%------------------------------------------------------------------------------%
\begin{question}
  Converta as coordenadas polares em coordenadas cartesianas $(x, y)$

  a) \quad \((r, \theta) = \left( 4, \frac{ \pi}{6} \right) \)

  b) \quad \((r, \theta) = \left( 3, \frac{3\pi}{4} \right) \)

  c) \quad \((r, \theta) = \left( 5, \pi            \right) \)
\end{question}

%------------------------------------------------------------------------------%
\begin{resolution}{Solução Item a}
  \[
    x      \;=\; r \cos(\theta)
    \pause \;=\; 4 \cos\left(\frac{\pi}{6}\right)
    \pause \;=\; 4 \frac{\sqrt{3}}{2}
    \pause \;=\; 2 \sqrt{3}
  \]

  \[
    \pause y \;=\; r \sin(\theta)
    \pause   \;=\; 4 \sin\left(\frac{\pi}{6}\right)
    \pause   \;=\; 4 \frac{1}{2}
    \pause   \;=\; 2
  \]

  \[
    \pause (x, y) = \left( 2 \sqrt{3}, 2 \right)
  \]
\end{resolution}

%------------------------------------------------------------------------------%
\begin{resolution}{Solução Item b}
  \[
    x      \;=\; r \cos(\theta)
    \pause \;=\; 3 \cos\left(\frac{3\pi}{4}\right)
    \pause \;=\; \frac{-3\sqrt{2}}{2}
  \]

  \[
    \pause y \;=\; r \sin(\theta)
    \pause   \;=\; 3 \sin\left(\frac{3\pi}{4}\right)
    \pause   \;=\; \frac{3\sqrt{2}}{2}
  \]

  \[
    \pause (x, y) = \left(-\dfrac{3\sqrt{2}}{2}, \dfrac{3\sqrt{2}}{2}\right)
  \]
\end{resolution}

%------------------------------------------------------------------------------%
\begin{resolution}{Solução Item c}
  \[
    x      \;=\; r \cos(\theta)
    \pause \;=\; 5 \cos\left(\pi\right)
    \pause \;=\; 5(-1)
    \pause \;=\; -5
  \]

  \[
    \pause y \;=\; r \sin(\theta)
    \pause   \;=\; 5 \sin\left(\pi\right)
    \pause   \;=\; 0
  \]

  \[
    \pause (x, y) = (-5, 0)
  \]
\end{resolution}

%------------------------------------------------------------------------------%